\section{The exact presentation and its automorphisms}
% BEGIN F1.Presentation
\begin{theorem}[The source-generator identification]\label{F1.Presentation}
Write
\[
 m-1=(e-1)(p-1)+r,\qquad e\geq1,\quad1\leq r\leq p-1,
\]
and put $d_i=p^e$ for $i\leq r$ and $d_i=p^{e-1}$ for $i>r$.
The Sylow subgroup in Theorem~\ref{F1.Construction} has the exact
presentation
\begin{align*}
 \langle s_1,\ldots,s_{p-1},\beta\mid{}
 &s_i^{d_i}=1,\quad\beta^p=1,\quad [s_i,s_j]=1,\\
 &[s_i,\beta]=s_{i+1}\ (1\leq i<p-1),\\
 &[s_{p-1},\beta]=\prod_{i=1}^{p-1}s_i^{-\binom pi}\rangle.
\end{align*}
The power and commutativity relators range over $1\leq i,j\leq p-1$.
The generator $s_i$ maps to $\pi^{i-1}$ in $A$ and has exactly
order $d_i$. Under this isomorphism, $\sigma$ fixes $s_1$ and
inverts $\beta$, while $\tau$ inverts $s_1$ and fixes $\beta$.
Thus the group, automorphisms and marking are precisely those of
\cite[Section 5 and Theorem 5.1]{HFLQ}.
\end{theorem}
\begin{proof}
The coordinate theorem of~\cite{UF} gives the exact additive orders
and commutativity of the proposed $s_i$. Our commutator convention
gives $[s_i,\beta]=\pi s_i=s_{i+1}$ for $i<p-1$.
The cyclotomic equation
\[
 p+\binom p2\pi+\cdots+\binom p{p-1}\pi^{p-2}
       +\pi^{p-1}=0
\]
gives the final relation. Thus all literal source relators hold.

Conversely let $\widehat P$ be the group defined by the displayed
presentation. Its subgroup $S=\gen{s_1,\ldots,s_{p-1}}$ is finite
abelian, with order at most
$\prod_i d_i=p^{er+(e-1)(p-1-r)}=p^{m-1}$.
The conjugation relators give $S^\beta\leq S$; since conjugation
is injective and $S$ is finite, this is equality. Hence $S$ is normal
and $\widehat P/S$ is cyclic of order at most $p$.
Consequently $|\widehat P|\leq p^m$. The homomorphism to the
constructed group $P$ is surjective: the displayed basis generates
$A$, and $\beta$ generates the quotient. Its image has order
$p^m$, so the homomorphism is an isomorphism. This proves the
generator-by-generator assertion, including modulus-one generators.

The formulas on $s_1,\beta$ are the definitions of $\sigma,\tau$.
These elements generate $P$, because successive commutators of
$s_1$ with $\beta$ give every $s_i$. The formulas therefore identify
the automorphisms, not only their effects on a proper subgroup.
\end{proof}
% END F1.Presentation

% BEGIN F1.Binomial
\paragraph{The auxiliary binomial formulas.}
Define $s_p=[s_{p-1},\beta]$ and $s_{p+1}=[s_p,\beta]$.
They are the additive classes $\pi^{p-1},\pi^p$; they are not
additional independent generators. The common conjugation identities
give, in multiplicative notation, for $1\leq k\leq p$,
\begin{align*}
 s_k^{\beta^{1-k}}&=\prod_{i=k}^{p+1}
           s_i^{\binom{p+1-k}{i-k}},\\
 s_k^\sigma&=
 \begin{cases}
 s_k^{\beta^{1-k}},&k\text{ odd},\\
 (s_k^{\beta^{1-k}})^{-1},&k\text{ even}.
 \end{cases}
\end{align*}
These are exactly Lemma 5.2(1),(2) of~\cite{HFLQ}.
The verified automorphisms in Theorem~\ref{F1.Construction} give
Lemma 5.3 directly. The source's coefficient-table calculation is
therefore unnecessary for this proof: the full relations are
preserved by actual automorphisms of the ring model and the exact
presentation isomorphism transports them back.
% END F1.Binomial

% BEGIN F1.Class
\paragraph{Nilpotency class and the Hughes hypothesis.}
The group $A$ is abelian, $P/A$ is cyclic, and
$[A,\beta]=\pi A$. Induction gives
\[
 \gamma_j(P)=\pi^{j-1}A\qquad(j\geq2).
\]
The successive lattice indices show that these groups are nontrivial
through $j=m-1$ and trivial at $j=m$. Thus $P$ has nilpotency
class $m-1$, the maximal-class property stated in the source.
The motivating condition $A=H_p(P)$ in its Section 5 is not a
condition needed for the constructed presentation. In fact, if
$m-1\leq p-1$, the coordinate theorem makes $A$ elementary
abelian. All elements outside $A$ have order $p$ as well, so every
element of $P$ has order dividing $p$. The Hughes subgroup, generated
by elements of order different from $p$, is then trivial. No general
classification of Hughes groups or groups of maximal class is used
in any of our arguments.
% END F1.Class
